We introduced a reinforcement learning approach for learning unmasking strategies in diffusion LLMs.
Our experiments demonstrated that the learned policies can match and sometimes exceed the performance of recently proposed sampling heuristics, paving the way for the automated discovery of scalable and robust sampling mechanisms. %

\paragraph{Limitations and future work} While we have shown that training sampling policies on LLaDA \citep{llada2025} allows us to match the performance of heuristics such as Fast-dLLM \citep{fastdllm2025}, a small performance gap remains when policies are trained on Dream \citep{ye2025dream}. Understanding which characteristics of Dream (e.g., it being initialized from an AR model) contribute to the diminished performance of RL policies is an important open question. 
One possible direction for future work is thus to train policies on a mixture of models, which might help bridge this gap, or to explore other lightweight ways of incorporating semantic information (e.g., by pairing confidences with the (un)embedding vectors of the corresponding tokens).
Another limitation of our approach is that it requires training a separate policy for each value of $\alpha$ (cf. \Cref{eq:reward}), and that this hyper-parameter sometimes does not yield as much fine-grained control over the policy's behavior as desired (c.f. \Cref{sec:exp-bl32}).
Another interesting avenue for future work is thus to investigate whether the accuracy-speed tradeoff could be controlled through some alternative mechanism, instead of explicitly having to set the strength of the computational penalty during training.
We also observe that the optimal policy temperature $\tau_{\pi}$ varies across generation settings (cf. \Cref{fig:combined-results}), suggesting the need to explore whether $\tau_{\pi}$ can be learned jointly with the policy, or whether alternative likelihood parameterizations (such as DPLS, see \autoref{sec:appendix-dpls}) offer greater robustness to temperature selection.
Beyond overcoming these limitations, promising future extensions include: (i) expanding the training data mixture to incorporate samples from multiple domains; (ii) extending our sampling policies to support not only unmasking but also remasking \citep{wang2025remasking}; and (iii) moving beyond the text domain to learn samplers for multimodal discrete diffusion models \citep{swerdlow2025unified}.
